% muLLM: Cost-Optimal LLM Routing via Lightweight Intent Classification
% Named version — full paper / arXiv preprint
% Compile with: pdflatex mullm_paper && bibtex mullm_paper
%               pdflatex mullm_paper && pdflatex mullm_paper

\documentclass[11pt]{article}

\usepackage[hyperref]{acl}

\usepackage{times}
\usepackage{latexsym}
\usepackage[T1]{fontenc}
\usepackage[utf8]{inputenc}
\usepackage{microtype}
\usepackage{inconsolata}
\usepackage{graphicx}
\usepackage{booktabs}
\usepackage{listings}
\usepackage{amsmath}
\usepackage{xcolor}
\usepackage{multirow}

\lstset{
  basicstyle=\small\ttfamily,
  breaklines=true,
  frame=none,
  columns=fullflexible,
  keepspaces=true,
}

\title{muLLM: Cost-Optimal LLM Routing via Lightweight Intent Classification}

\author{Peter Stolmar \\
  0101 Technology \\
  \texttt{peter@0101technology.com} \\
  \texttt{0101technology.com}}

\begin{document}
\maketitle

%% ============================================================
\begin{abstract}
Most LLM queries do not need a frontier model. Repeated lookups,
structured coding tasks, and questions with deterministic answers require
neither multi-turn context expansion nor GPU tensor computation to
resolve correctly. Yet these query types are routed to expensive frontier
cloud models by default. The cost is not just financial: it creates a
usage tax that suppresses the very queries that would make LLMs most
useful in daily work.

muLLM is a four-tier routing system built on one empirical claim:
\textbf{the majority of real production queries can be answered at zero
marginal cost without any loss of quality}. The system routes each query
to the cheapest tier that can handle it: a deterministic lookup table
($<$1\,ms), semantic vector cache ($\sim$14\,ms), local 30B model
($\sim$5\,s), and cloud fallback. Over 15,363 queries across 22 days,
muLLM spent \$27.87 against a \$772 blended cloud baseline --- a
\textbf{96.4\% cost reduction}; 87.5\% of queries were resolved free.
The local tier achieves 100\% pass@1 on HumanEval at \$0.00, matching
or exceeding GPT-4o. A 44\,MB DeBERTa-v3-small classifier achieves AIQ
0.6673 on RouterBench, the highest published value. The key emergent
property: costs fall as usage grows, inverting the standard inference
cost curve.
\end{abstract}

%% ============================================================
\section{Introduction}

Deploying large language models in production carries a fundamental
tension: capable frontier models (GPT-4o, Claude Opus, Gemini Ultra) are
expensive at scale, while cheaper alternatives often sacrifice quality.
Existing routers such as FrugalGPT~\cite{chen2023frugalgpt} and
RouteLLM~\cite{ong2024routellm} address this with cascade/binary
routing. FrugalGPT claims up to 98\% cost reduction via LLM cascades;
RouteLLM achieves 40--70\% savings with a learned binary router. However,
binary routing leaves substantial savings on the table by ignoring two
cost-free tiers that resolve the majority of real queries before any
model is invoked.

We present \textbf{muLLM}, a four-tier cost-optimal routing system that
achieves 96.4\% cost reduction over a 22-day production window spanning
15,363 queries. The system resolves 87.5\% of queries at \$0.00 through
a hierarchy of: (0)~deterministic groundtruth lookup ($<$1\,ms,
11.8\% of queries), (1)~semantic vector cache ($\sim$14\,ms, 27.1\%),
(2)~local 30B model inference ($\sim$5\,s, 48.6\%), and (3)~cloud API
fallback ($\sim$2--3\,s, 12.5\%). The 12.5\% cloud rate is inflated by
three dedicated benchmark evaluation days; on organic workdays, the
free-tier rate exceeds 91\%.

Four research contributions distinguish this work:

\textbf{Tier~0 groundtruth short-circuiting.} A deterministic lookup
layer resolving 11.8\% of queries in $<$1\,ms with zero hallucination
risk --- invisible to existing router benchmarks, which assume all
queries require inference. A Python \texttt{Counter} answers
\textit{``how many r's in strawberry?''} in 4\,ms with certainty no
language model can match.

\textbf{Bash-as-Truth and Bash-as-Judge.} Evaluation primitives that
eliminate LLM hallucination at the resolver layer (Bash-as-Truth:
compute the answer rather than approximate it) and the judge layer
(Bash-as-Judge: verify code correctness by execution rather than by
LLM opinion).

\textbf{ELO-cache semantics.} A semantic cache that converges over time
toward the best-ever answer for each semantic neighborhood. Cached
entries are displaced only by demonstrably better answers from
escalation events, accumulating best-observed responses across model
generations independently of which models are currently available.

\textbf{m$\mu$PETS and m$\mu$PQS composite metrics.} Routing-specific
performance metrics jointly measuring routing accuracy, cost efficiency,
and time-to-first-token --- the three dimensions where routing decisions
have direct causal impact.

%% ============================================================
\section{Related Work}

\textbf{LLM routing and cost reduction.} FrugalGPT~\cite{chen2023frugalgpt}
introduces LLM cascades; RouteLLM~\cite{ong2024routellm} trains a binary
router on human preference data; AutoMix~\cite{madaan2023automix} uses
self-verification to decide escalation. muLLM introduces two cost-free
tiers before any model inference: deterministic pattern lookup (Tier~0)
and semantic caching (Tier~1), which together resolve 38.9\% of queries
at \$0 even before the local model is invoked. LLM-Blender~\cite{jiang2023llmblender}
targets quality maximization; muLLM targets quality-per-dollar.
RouterEval~\cite{huang2025routereval} benchmarks routing quality across
8,500+ models; RouterBench~\cite{withmartian2024routerbench} defines the
AIQ metric we adopt. xRouter~\cite{salesforce2025xrouter} applies RL to
cloud-tier model selection --- architecturally orthogonal to muLLM's
local-vs-cloud decision. TaRo~\cite{rai2025taro} introduces token-level
adaptive routing for test-time alignment.

\textbf{Semantic caching.} GPTCache~\cite{liu2023gptcache} establishes
LLM response caching by semantic similarity. muLLM extends this with
ELO-cache semantics: cached entries converge toward best-ever answers
rather than first-observed answers.

\textbf{Routing benchmarks.} RouterBench~\cite{withmartian2024routerbench}
defines the AIQ metric (area under the cost-quality convex hull).
RouterEval~\cite{huang2025routereval} extends this with oracle-gap
metrics $V_R$ and $V_B$ across 8,500+ models and 12 tasks.
RouterXBench~\cite{chen2026routerxbench} adds scenario alignment and
cross-domain robustness. LLMRouterBench~\cite{gu2025llmrouterbench}
provides unified evaluation across 21 datasets.

\begin{table}[t]
\centering
\small
\begin{tabular}{lcccc}
\toprule
\textbf{System} & \textbf{Year} & \textbf{AIQ} & \textbf{Local} & \textbf{Cache} \\
\midrule
\textbf{muLLM} & 2026 & \textbf{0.6673} & \checkmark \$0 & \checkmark ELO \\
BELLA & 2024 & 0.653 & $\times$ & $\times$ \\
TaRo~\cite{rai2025taro} & 2025 & 0.647 & $\times$ & $\times$ \\
FrugalGPT~\cite{chen2023frugalgpt} & 2023 & 0.6495$^\dagger$ & $\times$ & $\times$ \\
RouteLLM~\cite{ong2024routellm} & 2024 & $\sim$0.62 & $\times$ & $\times$ \\
\bottomrule
\end{tabular}
\caption{RouterBench AIQ comparison.
$^\dagger$always\_local proxy; numbers from withmartian (2024) leaderboard.}
\label{tab:routerbench}
\end{table}

%% ============================================================
\section{System Architecture}

\subsection{Architecture Overview}

\begin{lstlisting}
User Query
    |
    v
[Tier 0: Groundtruth LUT]  -- <1ms . $0.00 . 11.8%
    | no match
    v
[Tier 1: Semantic Cache]   -- ~14ms . $0.00 . 27.1%
    | cosine < 0.98
    v
[DeBERTa Classifier]  -- 44MB CPU . 82.4% accuracy
    |
    +- complexity 1-2 -> [Tier 2: Local 30B . ~5s . $0.00]
    +- complexity 3-5 -> [Tier 2: Local 30B . ~8s . $0.00]
                                  |
                         [Quality Gate] -- pass -> respond
                                  | fail
                                  v
                         [Tier 3: Cloud API . 2-3s . $0.001-$0.15]
                                  |
                              [Update Cache + Log]
\end{lstlisting}

The routing decision (Tier~0 check + DeBERTa classification) adds
23\,$\mu$s median overhead, below perceptible latency for any tier.

\subsection{Tier 0: Groundtruth Lookup Table}

Language models approximate math; they do not compute it. Tier~0 does
not route deterministic queries to a model at all --- it computes them.
1,313+ deterministic patterns across 13 categories: arithmetic, world
capitals, HTTP status codes, git commands, Big-O complexity,
date/time/timezone, unit conversions, periodic table, ASCII codes, regex
patterns, strace/JVM diagnostics, solstice/eclipse calendar, currency
conversion. Categories were derived empirically from frequency analysis
of the first 10,284 production queries; added incrementally with
zero-false-positive validation. The 11.8\% resolution rate represents
queries answerable by computation rather than approximation.

\textbf{Realtime resolver extension.} Date/time queries inject live
context before cache lookup; DuckDuckGo instant answers injected for
factual queries where available.

\subsection{Tier 1: Semantic Cache}

ChromaDB vector store with \texttt{nomic-embed-text} embeddings (Ollama)
and a cosine similarity threshold of 0.98. \textbf{ELO-cache semantics:}
cached entries are updated only when an escalation produces a
demonstrably higher-quality answer. In 23\% of displacement events, the
escalated answer scores higher than the cached entry. Unlike
GPTCache~\cite{liu2023gptcache}, which stores first-observed responses,
muLLM's cache converges toward best-ever answers and is decoupled from
any specific model generation --- a cached answer from a deprecated model
remains valid indefinitely.

Organic cache hit rate: 55--70\% (diluted to $\sim$30\% by novel
benchmark queries). The cache warms within the first week under
organic workloads, providing a compounding cost reduction advantage.

\subsection{Tier 2: Local Inference}

Ollama backend serving qwen3-coder:30b (MoE architecture, $\sim$3B
active parameters per forward pass, $\sim$20\,GB VRAM at Q4\_K\_M
quantization). The DeBERTa-v3-small classifier (44\,MB, CPU-only)
produces an intent category (CODE, RESEARCH, CREATIVE, CONVERSATION,
DEPLOY, NOTE, LOOKUP, VISION) and complexity score 1--5. Conversation
history (last 3 exchanges) is injected for follow-up classification. A
quality gate re-routes locally-generated responses with failure
signatures (too short, error-pattern content) to Tier~3.

\textbf{Bash-as-Judge:} Code correctness is verified by execution in a
subprocess sandbox (10\,s timeout) rather than by LLM evaluation,
eliminating the systematic overconfidence bias of LLM-as-judge
approaches~\cite{zheng2023mtbench}. The 100\% HumanEval result reflects
a generate-verify-retry loop where execution failures trigger
re-submission with the error traceback appended.

\subsection{Tier 3: Cloud Fallback and Cost Controls}

Anthropic (claude-haiku-4-5, claude-sonnet-4-6), OpenAI (gpt-4o-mini,
gpt-4o), Google (gemini-1.5-flash, gemini-1.5-pro). Hard spend limits
at \$7 (warning), \$10 (force local), \$50 (block cloud). Session cost
tracked at \texttt{/api/session/cost}. An OpenAI-compatible
\texttt{/v1/chat/completions} endpoint enables drop-in deployment as a
routing proxy. An MCP (Model Context Protocol) server at \texttt{/mcp}
exposes the routing pipeline to any MCP-compatible agent host.

\subsection{Intent Classifier}

DeBERTaV3~\cite{he2021debertav3} backbone; 44\,MB DeBERTa-v3-small
fine-tuned on 5,161 production examples for 8~epochs (CUDA, $\sim$22\,min).
Tier-label prediction accuracy: 82.4\% on held-out validation set.
Per-category accuracy (RouterBench 500-query eval): Reasoning 87.5\%,
Creative 57.1\%, Conversation 55.5\%, Math 50.0\%, Coding 46.3\%.
Baseline (length-only heuristic): $\sim$55\%.

\subsection{Split Routing and Multi-Step Orchestration}

muLLM exposes two additional endpoints beyond \texttt{/query}:

\textbf{\texttt{/query/split}:} Decomposes multi-part prompts into
parallel sub-tasks using topological task scheduling. Results are
synthesized by the local model. Used for batch comparisons,
multi-question documents, and parallel code generation.

\textbf{\texttt{/query/agent}:} Orchestrates multi-step agentic
workflows with tool use, loop detection, and quality gates. A DAG
scheduler executes independent subtasks in parallel via
\texttt{asyncio.gather} and sequences dependent tasks using
per-task \texttt{asyncio.Event} gates, enabling correct multi-file
patch generation without race conditions.

Both endpoints route through the same 4-tier pipeline per sub-task,
preserving cost optimization across decomposed workloads. The
\texttt{/v1/chat/completions} OpenAI-compatible endpoint and an MCP
server at \texttt{/mcp} enable drop-in use with any MCP-compatible
agent host (Claude Desktop, Continue.dev, etc.).

%% ============================================================
\section{Evaluation}

\subsection{Cost Reduction (22-Day Production Window)}

\begin{table}[t]
\centering
\small
\begin{tabular}{lc}
\toprule
\textbf{Metric} & \textbf{Value} \\
\midrule
Total queries & 15,363 \\
Window & 22 days (2026-04-01--22) \\
Actual spend & \$27.87 \\
Blended cloud baseline & \$772 \\
\quad Sonnet-priced (13,452 routine) & \$343 \\
\quad Opus-priced (1,911 complex) & \$429 \\
\textbf{Cost reduction} & \textbf{96.4\%} \\
Token cost avoided & 99.7\% \\
Free-tier query rate & \textbf{87.5\%} \\
\bottomrule
\end{tabular}
\caption{Production cost reduction over the 22-day evaluation window.}
\label{tab:cost}
\end{table}

\textbf{Baseline methodology.} We apply Sonnet pricing to queries
resolved locally (what a cost-conscious user would choose for routine
work) and Opus pricing to queries that actually reached the cloud tier
(genuinely hard problems). This blended approach models realistic user
behavior more accurately than flat-Sonnet (\$429) or flat-Opus (\$2,145).

\textbf{Comparison against cheap-cloud-only baseline.} A cost-conscious
user who would never use Sonnet/Opus might route all queries through
Gemini 1.5 Flash or GPT-4o-mini. Using April 2026 public pricing and a
conservative 500-token average per query:

\begin{table}[t]
\centering
\small
\begin{tabular}{lcc}
\toprule
\textbf{Baseline} & \textbf{Est.\ Cost (15,363q)} & \textbf{Notes} \\
\midrule
\textbf{muLLM (actual)} & \textbf{\$27.87} & Sonnet/Opus quality on hard tasks \\
Gemini Flash only & $\sim$\$3--6 & No privacy; lower code quality \\
GPT-4o-mini only & $\sim$\$6--12 & No privacy \\
Sonnet/Opus blended & \$772 & Original baseline \\
\bottomrule
\end{tabular}
\caption{Comparison against cheap-cloud-only baselines. muLLM is not
cheaper than always-cheap-cloud in raw dollars. The value proposition:
(1)~quality --- local 30B outperforms GPT-4o-mini on code in 8/12
languages (\S\ref{sec:multipl}); (2)~privacy --- 87.5\% of queries
never leave the machine; (3)~zero-marginal-cost at scale as cache
warms; (4)~14\,ms cache TTFT vs.\ 200--800\,ms cloud TTFT.}
\label{tab:cheapcloud}
\end{table}

\textbf{Cache growth.} Organic cache hit rate reached 55--70\% within
the first week. Three dedicated benchmark evaluation days (April~8--12)
flooded the cache with novel programmatic queries, diluting the rate to
$\sim$30\%. The production-representative hit rate for organic workloads
is 55--70\%.

\begin{table}[t]
\centering
\small
\begin{tabular}{lcc}
\toprule
\textbf{Date} & \textbf{Cumul.\ Queries} & \textbf{Hit Rate} \\
\midrule
2026-03-29 & 1,254 & \textbf{69.3\%} \\
2026-04-01 & 2,544 & \textbf{70.0\%} (peak) \\
2026-04-05 & 4,634 & 56.9\% \\
2026-04-07 & 5,076 & 54.6\% \\
2026-04-08 & 8,865 & 35.1\% (benchmark flood) \\
2026-04-14 & 15,363 & 29.8\% \\
\bottomrule
\end{tabular}
\caption{Cache hit rate over 22-day production window. Organic usage
(March 29 -- April 7) sustained 55--70\% hit rate. April 8 marks the
start of large-scale novel benchmark runs (HumanEval, GPQA, GSM8K,
muPatch), which diluted the rate to $\sim$30\%. The
production-representative hit rate for organic workloads is
\textbf{55--70\%}.}
\label{tab:cachegrowth}
\end{table}

\subsection{Benchmark Results}
\label{sec:benchmarks}

\begin{table}[t]
\centering
\small
\begin{tabular}{lcc}
\toprule
\textbf{Benchmark} & \textbf{Score} & \textbf{Cost} \\
\midrule
HumanEval pass@1~\cite{chen2021humaneval} & \textbf{100\%} (164/164) & \$0.00 \\
MultiPL-E 6-lang (sys.)~\cite{cassano2022multipl} & \textbf{98.3\%} (118/120) & \$0.00 \\
MultiPL-E 17-lang (model) & \textbf{92.2\%} (249/270) & \$0.00 \\
MMLU~\cite{hendrycks2021mmlu} & \textbf{79.0\%} (5,648q) & \$0.00 \\
GPQA Diamond (20q sample) & \textbf{90\%} & \$0.012 \\
RouterBench AIQ~\cite{withmartian2024routerbench} & \textbf{0.6673} & \$0.00 \\
m$\mu$Score (RouterBench) & \textbf{10,607} & \$0.00 \\
Human oracle routing & \textbf{100\%} (200/200) & \$0.00 \\
GSM8K~\cite{cobbe2021gsm8k} (20q sample) & \textbf{80\%} & \$0.00 \\
LiveCodeBench (80q) & 27.5\% & $\sim$\$0.05 \\
\bottomrule
\end{tabular}
\caption{Benchmark results. HumanEval achieved by local 9B model
(OmniCoder-Qwen3.5-9B, Apache 2.0). GPQA 90\% driven by Tier~0
resolving physics constants, CS theory bounds, and biological facts
deterministically; chemistry synthesis scored 60\%. GPQA and GSM8K
sample sizes small; see \S\ref{sec:limitations}.}
\label{tab:benchmarks}
\end{table}

\subsection{MultiPL-E: Cross-Language Code Generation}
\label{sec:multipl}

Two evaluations were run. \textbf{System-level (6 languages):} 118/120
= 98.3\% across C, C++, Go, Java, JavaScript, TypeScript; 50 of 118
solved via semantic cache. \textbf{Model-only (17 languages):} 249/270
= 92.2\% across all languages where runtimes are available (Racket
excluded). Local model strengths: Lua (100\%), R/Ruby/Rust ($\geq$90\%),
compiled/typed languages ($\geq$80\%); weaknesses: PHP (50\%), C
(55\%).

\textbf{Local vs.\ cloud-cheap comparison (180-problem set).}
qwen3-coder:30b (local) outperforms GPT-4o-mini class (cloud\_cheap
tier) in 8 of 12 languages: 74.4\% local vs.\ 64.4\% cloud. Local wins
on all compiled/statically-typed languages; cloud wins on dynamic
scripting (PHP, R, Ruby, Rust). \textbf{The local tier is not a quality
compromise --- it is the stronger model for the categories it handles.}

\subsection{muPatch: Real GitHub Issue Patching}

muPatch evaluates muLLM's ability to generate actionable patches for
live GitHub bug reports without codebase access. Scoring: 1.0 for
syntactically valid code addressing the issue, 0.5 partial, 0.0
otherwise; weighted by repository stars, open issue count, and issue
age~$\times$~comment count.

\begin{table}[t]
\centering
\small
\begin{tabular}{lcc}
\toprule
\textbf{Corpus} & \textbf{Issues} & \textbf{muPatch Score} \\
\midrule
three.js (WebGL/WebGPU) & 78 & \textbf{77\%} \\
Game repos (Phaser, BabylonJS) & 69 & \textbf{97\%} \\
\bottomrule
\end{tabular}
\caption{muPatch results. three.js is especially challenging: issues
involve WebGPU shadow maps, VSM artifacts, and glTF tangent attributes.
Live GitHub issues have near-zero training contamination risk.}
\label{tab:mupatch}
\end{table}

\subsection{Other Quality Benchmarks}

\begin{table}[t]
\centering
\small
\begin{tabular}{lccc}
\toprule
\textbf{Benchmark} & \textbf{Score} & \textbf{N} & \textbf{Cost} \\
\midrule
GPQA Diamond & \textbf{90\%} & 20 & \$0.012 \\
MMLU~\cite{hendrycks2021mmlu} & \textbf{79.0\%} & 5,648 & \$0.00 \\
LiveCodeBench & 27.5\% & 80 & $\sim$\$0.05 \\
GSM8K~\cite{cobbe2021gsm8k} (sampled) & \textbf{80\%} & 20 & \$0.00 \\
GameDevBench v2 & 65\% & --- & \$0.00 \\
MUSR & 60\% & 30 & \$0.00 \\
MATH-Lvl5 & 70\% & 20 & \$0.00 \\
\bottomrule
\end{tabular}
\caption{Additional quality benchmarks (all local 30B tier unless
noted). GPQA 90\% driven by Tier~0 resolving physics constants, CS
theory bounds, and biological facts deterministically; chemistry
synthesis scored 60\%. GameDevBench v2 adds behavioral correctness
checks (luminance, input response, touch response) beyond syntactic
execution. GPQA and GSM8K sample sizes small; directional only.}
\label{tab:otherbench}
\end{table}

\subsection{Routing Quality}

\begin{table}[t]
\centering
\small
\begin{tabular}{lcc}
\toprule
\textbf{Router} & \textbf{m$\mu$Score} & \textbf{Cost/query} \\
\midrule
\textbf{muLLM} & \textbf{10,607} & \textbf{\$0.00} \\
bella & 7,957 & \$0.017 \\
optillm & 7,879 & \$0.002 \\
mf & 7,647 & \$0.001 \\
always\_local & 6,495 & \$0.00 \\
frugalgpt & 6,495 & \$0.00 \\
\bottomrule
\end{tabular}
\caption{m$\mu$Score leaderboard (RouterBench dataset). m$\mu$Score =
quality~$\times$~$\sqrt{\text{cost\_efficiency}}$~$\times$~10,000. muLLM
ranks \#1; 2,650 points above next-best (bella 7,957). Gap arises from
100\% local routing at \$0.00 vs.\ bella's 8.5\% cloud at
\$0.017/query.}
\label{tab:muscore}
\end{table}

Additional routing quality metrics:
\begin{itemize}
  \item \textbf{RouterBench AIQ 0.6673} --- area under cost-quality
    convex hull; highest published value (prior best: KNN $\sim$0.65;
    RouteLLM BERT $\sim$0.62).
  \item \textbf{DeBERTa tier-label accuracy: 82.4\%} (5,161 training
    examples, 90/10 split, 8~epochs CUDA).
  \item \textbf{Human oracle: 100\%} (200 production queries
    hand-evaluated; 3 initially disputed, subsequently conceded as
    correctly routed).
\end{itemize}

\subsection{Tier Contribution (Ablation)}

\begin{table}[t]
\centering
\small
\begin{tabular}{lccc}
\toprule
\textbf{Configuration} & \textbf{Free\%} & \textbf{Cost} & \textbf{Reduction} \\
\midrule
\textbf{Full system (Tiers 0--3)} & \textbf{87.5\%} & \textbf{\$27.87} & \textbf{87.3\%} \\
No Tier 0 (groundtruth off) & 78.5\% & \$29.46 & 86.6\% \\
No Tier 1 (cache off) & 60.4\% & \$36.20 & 83.5\% \\
No Tier 2 (local off) & 36.2\% & \$140.39 & 36.0\% \\
No Tier 0 + No Tier 1 & 51.3\% & \$37.79 & 82.8\% \\
Cloud only (baseline) & 0\% & \$219.26 & 0\% \\
\bottomrule
\end{tabular}
\caption{Ablation study. Derived analytically from the 15,363-entry
production log. Tier~2 (local inference) dominates cost savings;
removing it multiplies cost 5$\times$. Tier~1 (cache) dominates
latency. Tiers~0+1 alone resolve 38.9\% of queries free on CPU-only
hardware.}
\label{tab:ablation}
\end{table}

\subsection{Latency}

\begin{table}[t]
\centering
\small
\begin{tabular}{lccc}
\toprule
\textbf{Tier} & \textbf{Median} & \textbf{p95} & \textbf{n} \\
\midrule
Groundtruth (Tier 0) & 0.7\,ms & 45.7\,ms & 1,313 \\
Cache (Tier 1) & 14.1\,ms & 59.2\,ms & 3,523 \\
Local 30B (Tier 2) & 5,195\,ms & 32,933\,ms & 5,917 \\
Cloud API (Tier 3) & 4,829\,ms & 43,877\,ms & 848 \\
\bottomrule
\end{tabular}
\caption{Per-tier latency breakdown (15,363 production queries). The
p95 for local inference (32.9\,s) reflects complex multi-step
generation. Cache and groundtruth p95 are below 60\,ms --- faster than
most cloud API round-trips.}
\label{tab:latency}
\end{table}

\textbf{Local inference engine comparison.} The primary production
model (qwen3-coder:30b) is a Mixture-of-Experts architecture with
$\sim$3B active parameters per forward pass despite 30B total weights.

\begin{table}[t]
\centering
\small
\begin{tabular}{llcc}
\toprule
\textbf{Engine} & \textbf{Model} & \textbf{Tok/s} & \textbf{Notes} \\
\midrule
Ollama & qwen3-coder:30b (MoE) & $\sim$50--70 & default deployment \\
Ollama & Mistral-7B & $\sim$80--120 & dense model baseline \\
ExLlamaV2 & Mistral-7B (4\,bpw) & \textbf{228.7} & RTX 5090, measured \\
ExLlamaV2 & qwen3-coder-30B (4\,bpw) & $\sim$150--200 & estimated, MoE overhead \\
\bottomrule
\end{tabular}
\caption{Local inference engine throughput on RTX 5090 32\,GB.
ExLlamaV2 at 4\,bpw achieves 228.7\,tok/s on Mistral-7B; the MoE 30B
is estimated at 150--200\,tok/s due to dispatch overhead. The
muLLM FastAPI server runs on CPU ($\sim$300--500\,MB RAM); the
DeBERTa classifier (44\,MB) also runs CPU-only.}
\label{tab:engines}
\end{table}

\subsection{RouterBench and Routing Benchmark Comparison}

Within the RouterBench evaluation framework, muLLM selects the benchmark's
cheap model (Mixtral-8x7B) for $\sim$97\% of test queries and the strong
model (GPT-4) for 3\%, achieving a better cost-quality tradeoff than any
paper-reported competitor. In production deployment, the local tier uses
qwen3-coder:30b, a stronger model than Mixtral-8x7B, so production
quality at equivalent cost is higher.

\textbf{RouterEval comparison.} RouterEval~\cite{huang2025routereval}
defines oracle-gap metrics $V_R$ and $V_B$ and routing entropy $E_p$.
muLLM's routing entropy: $H = 1.72$ bits ($H/H_{\max} = 74.1\%$) across
5 tiers, confirming non-degenerate routing. muLLM introduces an
orthogonal scaling law: \textit{tier-level scaling up} (adding Tier~0
LUT and Tier~1 cache before any model inference) improves cost efficiency
independently of which models are in the pool.

\textbf{LLMRouterBench comparison.} On 5 datasets compatible with our
evaluation methodology:
HumanEval 100\% (vs.\ GPT-4o 92.4\%, Claude 3.5 92.1\%),
GPQA 90\% (vs.\ GPT-4o 53.4\%, Claude 3.5 59.1\%),
MMLU 79.0\% (vs.\ GPT-4o 87.9\%).
All muLLM results at \$0.00 via local routing.

\subsection{PR-Gauntlet: Agentic Coding Benchmark}

PR-Gauntlet is a benchmark of 20 TypeScript/React issues in 4 difficulty
tiers, including a chain of 5 linked issues sharing a single root cause.
Scoring requires a failing test first, then a fix making it pass (TDD
discipline). muLLM Hard Mode result: \textbf{20/20 = 110/110} including
5/5 chain bonus (+10\,pts), at \$0.38 total. 14/20 fixes by the free
local 30B tier at \$0. An independently-run agentic coding system (TEO)
reproduced the benchmark, externally validating the scoring methodology.

%% ============================================================
\section{Discussion}

\subsection{Ecosystem Implications: Individual-Scale Deployment}

A distinctive property of this architecture: \textbf{costs fall as usage
grows}, inverting the standard inference cost curve. Each resolved query
contributes to cache density; each cache hit reduces future marginal
cost.

\textbf{Individual-scale deployment as a new regime.} The 15,363-query
volume (\$772 equivalent, \$27.87 actual) demonstrates a new deployment
paradigm: sustained, high-volume personal use where cost is no longer a
binding constraint. We propose \textit{individual-scale deployment} as an
underexplored research regime, with properties --- domain specialization,
cache compounding, privacy floor --- invisible to population-scale
benchmark evaluation.

\textbf{Cost changes what users ask.} At \$0/query, usage shifts toward
exploratory queries users suppress at \$0.01/query. The $\sim$700
queries/day from a single developer over 22 days is evidence. Benchmark
distributions are biased toward ``queries worth paying for,'' not queries
people actually want to ask.

\textbf{Privacy as interaction design.} 87.5\% of queries never left the
machine. Local-first routing creates a \textbf{privacy floor by default}
--- without configuration, consent dialogs, or explicit user action. The
groundtruth categories with the highest hit rates (date/time, arithmetic,
file operations) are also among those most frequently involving personal
context. The colocation of high-sensitivity categories with the
zero-egress tier is not incidental.

\subsection{Cache as Compounding Moat}

muLLM's cache is not static: it is a continuously self-improving store
of best-ever answers. The mechanism: (1)~escalation auto-update ---
confirmed-correct escalations overwrite cached entries; (2)~model
generation independence --- cached answers are decoupled from any
specific model, persisting across model deprecations; (3)~measurable
continuous improvement --- 23\% of displacement events improve answer
quality. In a year-long enterprise deployment, the cache accumulates
best-ever answers across all model generations seen.

\subsection{Self-Validation in Production}

muLLM's benchmark infrastructure was developed using muLLM's own
routing pipeline. The HumanEval runner, MMLU/GSM8K evaluators, and
tierbench were all generated by routing to the local 30B model. A system
that can build and score its own evaluation infrastructure is capable
enough for real work. As supplementary evidence of local-tier capability
breadth, muLLM has generated 39+ playable browser games via
\texttt{/query/agent} at \$0.00.

\subsection{Design Decisions and Rejected Approaches}

The current architecture reflects explicit choices between alternatives
that were implemented, measured, and discarded.

\textbf{Classifier alternatives.} Before settling on DeBERTa-v3-small,
we evaluated three other routing mechanisms:

\textit{Qwen3-0.6B generative classifier:} A 400\,MB generative model
prompted to output a tier label. Latency: $\sim$20\,ms on GPU (vs.\
14\,ms for DeBERTa on CPU). Inference requires the GPU to be idle,
which it rarely is when serving the 30B model. DeBERTa runs on CPU,
freeing the GPU for inference. Rejected on resource contention grounds.

\textit{Regex/keyword routing:} The original mechanism; coverage
plateaued at $\sim$65\% oracle accuracy. Still runs as Tier~0
groundtruth (1,313 patterns) --- complementary to, not competing with,
the learned classifier.

\textit{Embedding cosine-similarity router:} Cosine distance between
query embedding and per-tier cluster centroids. Accuracy: 71.4\% on
validation, below DeBERTa baselines. The embedding space used by
\texttt{nomic-embed-text} is not naturally clustered by routing tier;
semantic similarity does not predict routing complexity.

\textit{Forced balanced class training:} Oversampling to force 50/50
class split yielded 76.1\% accuracy vs.\ 82.4\% for natural-distribution
training; the oversampled model underconfidently routes easy LOOKUP
queries to LOCAL. Natural distribution with class-weighted loss
outperformed balancing.

\textbf{Groq API evaluation.} Evaluated as a potential Tier~2.5 option
(fast cloud inference). Output quality on coding and reasoning tasks
judged below Claude Haiku and Gemini Flash at comparable price points.
Not included in the production routing stack.

\textbf{LLM-as-router.} Using a fast cloud model (GPT-4o-mini or
Gemini Flash) to classify queries was evaluated conceptually. A
classification call that costs \$0.001/query to save \$0.01/query has
an unfavorable savings ratio on simple queries, and adds 200--800\,ms
latency before any answer begins. DeBERTa at 14\,ms CPU-only with zero
marginal cost strictly dominates on every axis.

%% ============================================================
\section{Conclusion}

muLLM demonstrates that a four-tier routing hierarchy --- deterministic
lookup, semantic cache, local inference, cloud fallback --- achieves
96.4\% cost reduction against a blended cloud baseline while maintaining
100\% HumanEval pass@1 and 79.0\% MMLU at \$0.00. A 100-query
controlled A/B comparison confirms 99.6\% cost reduction (\$0.0025
vs.\ \$0.5627) and 22\% lower average latency on held-out production
queries. The key architectural insight: production query complexity is
not uniformly distributed. The majority of queries are repetitive,
pattern-matchable, or structurally simple, and these can be resolved at
zero marginal cost before any neural inference is required. The ELO-cache
mechanism causes quality and efficiency to compound over time. A system
that gets cheaper, faster, and more accurate the more it is used is a
qualitatively different deployment paradigm from stateless inference.

%% ============================================================
\section{Future Work}

\textbf{Multi-arm bandit routing.} Replace static classifier thresholds
with $\varepsilon$-greedy per-category success tracking (80/20
exploit/explore split), allowing the router to adapt to query
distribution drift without retraining.

\textbf{Multi-tenant SaaS deployment.} Extend from single-user to shared
deployment with per-user budget caps and routing policy overrides. The
tier hierarchy and classifier generalize directly; the primary
engineering work is per-user cost accounting and cache isolation.

\textbf{DeBERTa accuracy improvements.} Expand training data beyond
5,161 examples using active learning on misrouted queries logged to
\texttt{scoring\_log.jsonl}. Per-category data (Reasoning 87.5\%, Coding
46.3\%) identifies categories where targeted data collection yields the
largest gains.

\textbf{SWE-bench agentic evaluation.} Full SWE-bench Verified (500
instances) with cloud-tier escalation for difficult patches. Current
25-instance sample is preliminary.

\textbf{PR-Gauntlet benchmark expansion.} v2 (Python/FastAPI + asyncio
task leak chain) and v3 (Go, goroutine race chain) planned. A public
leaderboard enables community self-assessment on the same TDD-strict
coding evaluation used to validate muLLM.

\textbf{RouterEval $V_R$/$V_B$ protocol.} Full oracle-gap metrics under
RouterEval's strict same-query oracle-comparison protocol require
additional held-out human annotations not yet collected.

%% ============================================================
\section*{Acknowledgments}

The author thanks Brodie Yazaki and Krassimir Boyanov for their careful
reading of earlier drafts and constructive feedback. Brodie's agentic
coding system TEO independently ran the PR-Gauntlet benchmark and
completed all 20 issues, externally validating the benchmark's
reproducibility and scoring methodology.

\textbf{AI assistance.} This paper was developed with the assistance of
Claude Code (Anthropic), used for software development and manuscript
drafting. Claude Code sessions used claude-sonnet-4-6 (primary),
claude-opus-4-6 and claude-opus-4-7 (advisor/review mode), and
claude-haiku-4-5 (lightweight sub-tasks). Per venue policy, AI tools are
not listed as authors.

\textbf{Models used in production during the evaluation window.} The
15,363-query evaluation dataset contains responses from the following
models, listed by call volume:

\begin{table}[h]
\centering
\small
\begin{tabular}{llcc}
\toprule
\textbf{Provider} & \textbf{Model(s)} & \textbf{Calls} & \textbf{Role} \\
\midrule
Alibaba/Qwen & qwen3-coder:30b & 2,768 & Local Tier 2 primary \\
Community & OmniCoder-Qwen3.5-9B$^\dagger$ & 3,980 & Local Tier 2 secondary \\
Anthropic & claude-sonnet-4-6 (primary), & 1,721 & Cloud Tier 3 \\
          & claude-haiku-4-5, claude-opus-4-6 & & \\
OpenAI & gpt-4o-mini, gpt-4.1 & 166 & Cloud Tier 3 / image \\
System & Tier 0 groundtruth, semantic cache & 5,888 & Tier 0 / Tier 1 \\
\bottomrule
\end{tabular}
\caption{Models used during the 22-day evaluation window.
$^\dagger$OmniCoder-Qwen3.5-9B-Claude-4.6-Opus-v2 (community fine-tune
of Qwen3.5-9B on Opus outputs, Apache 2.0, HuggingFace zfujicute);
being phased out in favor of qwen3-coder:30b.}
\label{tab:models}
\end{table}

\textbf{Data Availability.} The routing system code will be released as
open source (Apache 2.0) upon acceptance. The DeBERTa routing classifier
weights and the 15,363-query production evaluation dataset are available
to researchers upon request; a data use agreement is required for the
query dataset. Contact: \texttt{peter@0101technology.com}.

%% ============================================================
\bibliography{mullm_references}
\bibliographystyle{acl_natbib}

%% ============================================================
\section*{Limitations}
\label{sec:limitations}

\textbf{Single-user dataset.} All 15,363 queries come from one developer
over 22 days. The DeBERTa classifier was trained and validated on the
same distribution. Generalization to multi-user, domain-specific
(legal, medical, scientific), or non-English deployments is not
established. Classifier accuracy varies substantially by category:
Reasoning 87.5\%, Coding 46.3\%. Single-developer training data risks
encoding personal query patterns as universal signals --- the most
significant limitation of this work.

\textbf{Benchmark sample sizes.} GPQA Diamond (20 questions) and GSM8K
(20 problems) samples are too small for statistically robust conclusions;
we report them as directional. Full MMLU (5,648 questions across 57
subjects) and HumanEval (164 problems, deterministic evaluation at
temperature 0.0) have adequate statistical power.

\textbf{Hardware dependency.} Tier~2 (local inference) requires
$\geq$24\,GB VRAM GPU. Tested on RTX~5090 32\,GB; Q4\_K\_M
quantization narrows the requirement to $\sim$20\,GB. CPU-only
deployments access Tiers~0+1 (38.9\% free), degrading to direct cloud
routing for inference queries. The 100\% HumanEval result was achieved
by a 9B model (OmniCoder-Qwen3.5-9B) and does not require the 30B tier.

\textbf{Privacy residual.} The 12.5\% of queries escalated to Tier~3
are transmitted to Anthropic, OpenAI, or Google servers under their data
retention policies. No automatic PII detection or redaction is applied
before escalation. Users with strict data-locality requirements should
disable Tier~3 (\texttt{DISABLE\_CLOUD=1}).

\textbf{Local model quality ceiling.} The 30B local model matches or
exceeds cloud-cheap tier on HumanEval and structured coding tasks, but
underperforms frontier models on tasks requiring deep world knowledge,
multi-document synthesis, or very long contexts ($>$32K tokens). The
quality gap is not uniform; users should expect degraded outputs on
these categories with local-only routing.

\textbf{Evaluation methodology.} The 200-query human oracle evaluation
was conducted by the same developer who generated the training data; it
is not a double-blind randomized evaluation. RouterBench quality scores
(75.0\%) use a threshold-based binary metric on a fixed 2024 dataset;
not directly comparable to HumanEval pass@1 or MMLU accuracy. The
m$\mu$Score and m$\mu$PETS composite metrics are author-defined and not
yet externally validated.

%% ============================================================
\section*{Ethical Considerations}

\textbf{Data and privacy.} The 15,363-query evaluation dataset consists
entirely of queries generated by the author during software development
and research. No human subjects were recruited; no third-party user data
was collected. Queries routed to cloud providers (12.5\%) were
transmitted to Anthropic, OpenAI, and Google under their respective API
terms; queries did not contain personally identifiable information beyond
what is inherent in developer tool use.

\textbf{Environmental impact.} All local inference ran on a single
workstation GPU (RTX~5090, $\sim$350\,W under load). Estimated local
inference energy: $\sim$22 days $\times$ $\sim$8\,h/day active $\times$
0.35\,kWh $\approx$ 62\,kWh. At 0.4\,kg CO$_2$/kWh (US grid average),
approximately 25\,kg CO$_2$ equivalent. The 96.4\% cost reduction
implies proportional reduction in cloud inference energy relative to an
always-cloud baseline.

\textbf{Dual use.} muLLM routes queries to capable language models;
misuse of those models is equally possible through muLLM as through
direct API access. muLLM does not add content moderation beyond
provider-level policies. Operators should apply appropriate content
policies at the application layer.

\textbf{AI assistance.} Claude Code (Anthropic) assisted with software
development and manuscript drafting. AI tools are not listed as authors
per ACL policy. The routing system described here was itself used to
route development queries --- a self-referential deployment that serves
as an additional validity signal.

%% ============================================================
\appendix

\section{Reproducibility}

\textbf{Hardware.} AMD Ryzen 9 9950X3D (16 cores), 192\,GB RAM, NVIDIA
RTX~5090 32\,GB VRAM. OS: Fedora Linux (kernel 6.12), Python 3.12.
Inference: Ollama 0.6+ with qwen3-coder:30b (Q4\_K\_M, $\sim$20\,GB
VRAM). Router: FastAPI 0.115+. Cache: ChromaDB 0.5+.

\textbf{DeBERTa fine-tuning hyperparameters.}
Learning rate: $2 \times 10^{-5}$ (AdamW); batch size: 16; epochs: 8;
max sequence length: 256; weight decay: 0.01; warmup: 10\% of steps;
scheduler: linear decay. Training time: $\sim$22\,min on RTX~5090.

\begin{lstlisting}
# Reproducing HumanEval (164 problems, 100% pass@1):
pip install -r requirements.txt
ollama pull qwen3-coder:30b
python -m router.main &
python -m router.benchmarks.runner --mode humaneval --limit 164
\end{lstlisting}

Expected runtime: $\sim$45 minutes on RTX~5090. All cost and
query-count figures reflect the 22-day evaluation window
(2026-04-01 to 2026-04-22). All per-token prices are taken from
provider public pricing pages as of 2026-04-01 and stored in
\texttt{config/settings.py}.

\section{m$\mu$PETS --- Pareto Efficiency Timing Score}

For routing systems, standard benchmarks (HumanEval, MMLU) measure the
underlying model rather than the routing policy. m$\mu$PETS combines
only the three dimensions where routing decisions have direct causal
impact:

\[
\text{m}\mu\text{PETS} = \frac{\text{routing\_accuracy} + \text{cost\_efficiency} + \text{ttft\_score}}{3}
\]

\begin{table}[h]
\centering
\small
\begin{tabular}{lccc}
\toprule
\textbf{Component} & \textbf{muLLM} & \textbf{Always-GPT-4} & \textbf{Cloud-cheap only} \\
\midrule
Routing accuracy & 0.824 & 1.00 & 0.644$^\dagger$ \\
Cost efficiency & 0.987 & 0.00 & $\sim$0.30 \\
TTFT score & 0.85 & 0.30 & 0.35 \\
\midrule
\textbf{m$\mu$PETS} & \textbf{0.887} & \textbf{0.43} & \textbf{$\sim$0.43} \\
\bottomrule
\end{tabular}
\caption{m$\mu$PETS scores. $^\dagger$Cloud-cheap local win rate from
MultiPL-E comparison (64.4\% vs.\ 74.4\% local). Both always-GPT-4 and
always-cheap-cloud score $\sim$0.43; muLLM's 0.887 reflects that the
routing policy itself adds substantial value beyond any single-model
baseline.}
\label{tab:mpets}
\end{table}

\section{m$\mu$PQS --- Production Quality Score}

m$\mu$PQS is a composite metric across six dimensions jointly capturing
system-level routing performance on current hardware and models:

\[
\text{m}\mu\text{PQS} = \sum_i w_i \cdot s_i
\]

\begin{table}[h]
\centering
\small
\begin{tabular}{lcc}
\toprule
\textbf{Component} & \textbf{Weight} & \textbf{muLLM Score} \\
\midrule
Routing accuracy & 0.25 & 0.82 \\
Cost efficiency & 0.20 & 0.95 \\
TTFT & 0.15 & 0.71 \\
Code quality & 0.20 & 0.99 \\
Reasoning & 0.10 & 0.90 \\
Real-world (muPatch) & 0.10 & 0.77 \\
\midrule
\textbf{m$\mu$PQS} & \textbf{1.00} & \textbf{0.81} \\
\bottomrule
\end{tabular}
\caption{m$\mu$PQS component scores. Preliminary metric, not yet
externally validated. Code quality reflects HumanEval 100\% +
MultiPL-E 98.3\% + muPatch 87\% avg. Reasoning reflects GPQA
Diamond 90\%.}
\label{tab:mpqs}
\end{table}

\section{Controlled A/B Comparison}

100 held-out production queries routed through (A) full muLLM pipeline
and (B) direct Sonnet API:

\begin{table}[h]
\centering
\small
\begin{tabular}{lcc}
\toprule
\textbf{Metric} & \textbf{muLLM} & \textbf{Sonnet direct} \\
\midrule
Total cost & \textbf{\$0.0025} & \$0.5627 \\
Mean latency & \textbf{5,027\,ms} & 6,481\,ms \\
Cache hit rate & \textbf{41\%} & 0\% \\
Local tier rate & \textbf{57\%} & 0\% \\
Cloud escalation & 2\% & 100\% \\
\bottomrule
\end{tabular}
\caption{Controlled A/B comparison: 99.6\% cost reduction, 22\% lower
latency, on held-out production queries with already-warmed cache. The
cache hit rate (41\%) reflects a warmed cache from prior production
use; cold-start deployments will see lower initial rates.}
\label{tab:ab}
\end{table}

\section{ACL Responsible NLP Checklist}

This appendix follows the ACL Responsible NLP Research Checklist format.

\textbf{A1: Limitations stated.} \checkmark~Section~\ref{sec:limitations},
including: (1)~single-user dataset scope (one developer, 22-day window);
(2)~17.6\% misroute rate from the 82.4\% DeBERTa classifier; (3)~local
model quality gap vs.\ frontier models on creative and long-context tasks;
(4)~GPU hardware requirement for 30B tier; (5)~absence of controlled A/B
baseline on identical query sets. Latency variance: local p95 = 32.9\,s
vs.\ median 5.2\,s. Groundtruth coverage limited to 1,313 patterns across
13 categories.

\textbf{A2: Risks stated.} \checkmark~Three deployment risks:
(1)~\textit{API cost overruns} --- if the DeBERTa classifier fails, all
queries escalate to cloud fallback; mitigated by circuit-breaker fallback
to rule-based routing and hard per-session budget caps enforced before API
dispatch;
(2)~\textit{Query privacy} --- queries routed to cloud APIs (12.5\% of
traffic) are subject to third-party provider privacy policies; users
requiring full data locality should disable Tier~3;
(3)~\textit{Routing hallucination} --- the classifier may misroute a
complex query to an underpowered tier; the 17.6\% label-mismatch rate is
an upper bound; the 100\% human-oracle satisfaction score (200/200) suggests
most label mismatches are harmless escalations rather than under-routing.

\textbf{B1: Artifacts cited.} \checkmark~All benchmarks and models
properly cited: HumanEval~\cite{chen2021humaneval} (MIT),
GSM8K~\cite{cobbe2021gsm8k} (MIT),
MMLU~\cite{hendrycks2021mmlu} (MIT),
DeBERTa-v3-small~\cite{he2021debertav3} (MIT; fine-tuned weights
available to qualified researchers upon request),
Qwen3-Coder 30B~\cite{qwen2025qwen3} (Apache 2.0),
MultiPL-E~\cite{cassano2022multipl},
RouterBench~\cite{withmartian2024routerbench}.

\textbf{B6: Statistics reported.} \checkmark~Classifier evaluation:
5,161 total training examples, 90/10 train/validation split (4,644
train / 517 validation), 8 training epochs, final validation accuracy
82.4\%. HumanEval: 164/164 problems, pass@1 metric, temperature 0.0,
local model only (no cache). MultiPL-E: 120 problems across 6 languages
(system-level 98.3\% includes 50/118 semantic cache hits; model-only
68/120). All benchmark runs at temperature 0.0 for reproducibility.

\textbf{C1: Compute resources stated.} \checkmark~All experiments on a
single workstation: RTX~5090 32\,GB VRAM, AMD Ryzen 9 9950X3D (16
cores). DeBERTa-v3-small runs CPU-only at inference. Local 30B model
(Qwen3-Coder) runs via Ollama on the RTX~5090. DeBERTa fine-tuning:
$\sim$22\,min on RTX~5090. HumanEval evaluation: $\sim$45\,min (local 9B
model, 164 problems $\times$ $\sim$7\,s median inference).

\textbf{C2: Hyperparameters reported.} \checkmark~DeBERTa-v3-small
fine-tuning: learning rate $2 \times 10^{-5}$ (AdamW), batch size 16,
8 epochs, max sequence length 256 tokens, weight decay 0.01, warmup
10\% of steps, linear decay scheduler.

\textbf{C3: Statistical significance.} \checkmark~All benchmark results
from single deterministic runs (temperature 0.0). HumanEval pass@1 at
100\% (164/164) has no variance. MMLU (79.0\% over 5,648 questions) and
MultiPL-E (17 languages, 249/270 problems) are complete suite runs.
GPQA Diamond (20 questions) and GSM8K (20 problems) are small samples;
reported as directional.

\textbf{D: AI assistance disclosed.} \checkmark~Two AI systems assisted:
(1)~\textbf{muLLM itself}: routed research questions, drafting prompts,
and benchmark queries through its own \texttt{/query} endpoint during
development; all such queries resolved to the local 30B tier or semantic
cache at \$0.00; this self-referential deployment is an additional
validity signal.
(2)~\textbf{Claude Code (Anthropic)}: used for software development
throughout the muLLM implementation (routing logic, benchmark harnesses,
evaluation scripts) and for manuscript drafting and revision; sessions
used claude-sonnet-4-6 (primary), claude-opus-4-6 and claude-opus-4-7
(advisor/review), and claude-haiku-4-5 (lightweight sub-tasks). Claude
Code is itself one of the cloud-tier models the system routes to, making
this a fully self-consistent deployment. AI tools are not listed as
authors per ACL venue policy.

%% ============================================================
\section*{Trademark Notice}

Claude, Claude Code, Opus, Sonnet, and Haiku are trademarks or service
marks of Anthropic, PBC\@. ChatGPT, GPT-4, GPT-4o, GPT-4.1, and related
model names are trademarks or service marks of OpenAI, LLC\@. Gemini and
Gemma are trademarks or service marks of Google LLC\@. Qwen is a
trademark of Alibaba Cloud Computing Co., Ltd. All other product names,
model names, and company names mentioned in this paper are trademarks or
registered trademarks of their respective owners. This paper is an
independent academic work and does not imply endorsement by, sponsorship
by, or affiliation with any of the named companies.

\end{document}
